\section{Preliminaries}
\label{sec:preliminaries}

\paragraph{Notation.}
Let $\lambda$ denote the security parameter and let $\negl(\lambda)$ be any positive function negligible in
$\lambda$.  We write PPT for probabilistic polynomial time.  For integers $a<b$, write
$[a,b]=\{a,a+1,\ldots,b\}$ and use $[b]$ as shorthand for $[1,b]$.  Vectors are written in bold lower-case
letters and matrices in bold upper-case letters.

We work over the cyclotomic ring $\R=\Z[X]/(X^N+1)$ for power-of-two $N$ and its quotient
$\Rq=\R/q\R$, identified with polynomials of degree $<N$ with coefficients in $\Z_q$.  For $a\in\Z$, let
$\langle a\rangle_q\in\{0,\ldots,q-1\}$ denote reduction modulo $q$.  When a field element represents a
signed value, we use the centered representative.  For $x\in\Rq$, the notation $\|x\|_\infty\le B$ means that
every coefficient of $x$ lies in $[-B,B]$; the same convention is applied component-wise to vectors over
$\Rq$.

The full signed message is written
\[
  \vecmu=\vecm\Vert \veck,
\]
where $\vecm$ encodes holder attributes and $\veck$ is the secret PRF key used for rate limiting.  The public
coefficient bound for message and hash randomness is denoted $\BoundMsg$.

\subsection{Commitment schemes and Ajtai commitments}
\label{subsec:prelim-commit}
\label{subsec:commitments}

\begin{definition}[Commitment scheme]\label{def:commitment-scheme}
A non-interactive commitment scheme is a tuple of PPT algorithms
\[
  \ComScheme=(\ComSetup,\Commit,\ComVerify)
\]
with the following syntax.  The setup algorithm outputs public parameters $\ComPublicParam$; the commit
algorithm on input a message $m$ outputs a commitment $c$ and an opening $d$; and
$\ComVerify(c,m,d)$ deterministically checks validity of the opening.  Correctness, computational binding,
and hiding are defined in the usual way.
\end{definition}

\subsubsection{Ajtai's commitment scheme.}
We use the standard bounded linear commitment over $\Rq$, matching the compact Ajtai-style commitment
of~\cite[Definition~13]{EPRINT:AttCraKoh21}.  Fix an input length $\ell_{\mathsf{in}}$, randomness dimension
$\ell_r$, output length $\ell_{\ComScheme}$, and coefficient bound $\BoundMsg$.  Setup samples
\[
  \mACom \getsunif \Rq^{\ell_{\ComScheme}\times \ell_{\mathsf{in}}}
\]
and outputs $\ComPublicParam=(\mACom,\ell_r,\ell_{\ComScheme},\BoundMsg)$.  In the blind-signature and
ARC protocols, the committed vector has the form
\[
  \vecw=[\vecmu\Vert r_0^H\Vert r_1^H\Vert \bar r]^\top,
\]
where $r_0^H,r_1^H$ are holder-side rational-hash randomnesses and $\bar r$ is commitment-only randomness.
The holder computes
\[
  \vecc=\mACom\vecw.
\]
Equivalently, if
\[
  \mACom=[A_\mu\mid A_{r_0}\mid A_{r_1}\mid A_{\bar r}],
\]
then
\[
  \vecc=A_\mu\vecmu+A_{r_0}r_0^H+A_{r_1}r_1^H+A_{\bar r}\bar r.
\]
The verification algorithm accepts an opening iff the committed vector is coefficient-wise bounded by
$\BoundMsg$ and the displayed linear identity holds.

\begin{definition}[Module-SIS and Module-ISIS over $\Rq$]\label{def:sis-isis}
Fix integers $n,m\ge 1$, modulus $q\ge 2$, and bound $\beta>0$.
\begin{itemize}[leftmargin=1.25em]
  \item $\mathsf{M\text{-}SIS}_\infty(n,m,q,\beta)$: given $A\in\Rq^{n\times m}$, find non-zero
  $s\in R^m$ such that $As=0$ in $\Rq^n$ and $\|s\|_\infty\le \beta$.
  \item $\mathsf{M\text{-}ISIS}_\infty(n,m,q,\beta)$: given $(A,t)\in\Rq^{n\times m}\times\Rq^n$, find
  $s\in R^m$ such that $As=t$ in $\Rq^n$ and $\|s\|_\infty\le\beta$.
\end{itemize}
\end{definition}

\begin{lemma}[Binding of the Ajtai commitment]\label{lem:ajtai-binding}
Let $\Adv$ be any PPT adversary that opens an Ajtai commitment to two distinct bounded vectors with
probability $p$.  Then there is a PPT adversary that solves the corresponding
$\mathsf{M\text{-}SIS}_\infty$ instance with essentially the same probability and runtime.
\end{lemma}

\begin{proof}
Two valid openings of the same commitment give two distinct bounded vectors $\vecw,\vecw'$ such that
$\mACom(\vecw-\vecw')=0$ in $\Rq^{\ell_{\ComScheme}}$.  The difference is non-zero and coefficient-wise
bounded by $2\BoundMsg$, hence it is an SIS solution for the commitment matrix.
\end{proof}

\begin{lemma}[Statistical hiding of the Ajtai commitment]\label{lem:ajtai-hiding}
Let $B_r$ be the coefficient bound of the commitment-randomness alphabet.  If
\[
  N\ell_r\log_2(2B_r+1)\ge N\ell_{\ComScheme}\log_2 q+2\lambda,
\]
then the commitment distribution is statistically $2^{-\lambda}$-close to a distribution independent of the
committed message.  In the instantiation above, $B_r=\BoundMsg$ is sufficient.
\end{lemma}

\begin{proof}
This is the direct leftover-hash-lemma argument for the bounded linear Ajtai commitment, specialized to
$\Rq$.  The commitment randomness contributes $N\ell_r\log_2(2B_r+1)$ bits of min-entropy and the
commitment output space has size $q^{N\ell_{\ComScheme}}$.
\end{proof}

\subsection{Lattice trapdoors and rational hashing}
\label{subsec:prelim-vsis}
We recall only the definitions needed to state the construction.  Additional lattice-assumption syntax is
deferred to Appendix~\ref{app:assumptions lattices}.

\begin{definition}[Lattice trapdoor suite]\label{def:trapdoor-suite}
A lattice trapdoor suite over $\Rq$ consists of PPT algorithms $(\Setup,\TrapGen,\SampPre)$ such that:
\begin{itemize}[leftmargin=1.25em]
  \item $\Setup$ generates admissible public parameters $(n,m,\BoundSig)$.
  \item $(\mA,\trapdoor)\gets\TrapGen$ outputs a public matrix $\mA\in\Rq^{n\times m}$ and a trapdoor,
  with $\mA$ statistically close to uniform over $\Rq^{n\times m}$.
  \item On input $(\trapdoor,T)$, $\SampPre$ outputs $\vecu\in R^m$ such that $\mA\vecu=T$ in $\Rq^n$
  and $\|\vecu\|_\infty\le\BoundSig$, except with negligible failure probability.
\end{itemize}
\end{definition}

\begin{definition}[Trapdoor rational hash]\label{def:trapdoor-rational-hash}
A trapdoor rational hash scheme over $\Rq$ is a tuple
\[
  \TDRH=(\Setup,\TrapGen,\Eval,\SampPre)
\]
consisting of the lattice trapdoor suite above together with a public evaluation algorithm.  Its setup also
outputs a public hash description $B$, a randomness space $\mathcal X$, and an admissibility error bound
$\delta$.  For fixed $B$, the algorithm $\Eval(B,\mu,\chi)$ takes a message $\mu$ and witness randomness
$\chi\in\mathcal X$ and outputs either a target $T\in\Rq^n$ or $\bot$.
\end{definition}

\subsubsection{The BB-tran rational hash.}
We instantiate the rational hash with the BB-tran form of~\cite[Table~1]{EPRINT:DKLW25}.  The generic
DKLW expression is
\[
  h_{\mu,\chi}(\tilde b^\top)=(u^\top,x_0^\top)\tilde b_0+\frac{1}{\tilde b_1-x_1}.
\]
In the notation used in this paper, the public description is
$B=(B_0,B_1,B_2,B_3)$, where $B_1$ and $B_2$ are public linear maps into $\Rq$.  For message $\vecmu$ and randomness
$(r_0,r_1)$, define
\[
  L_B(\vecmu,r_0):=B_0+B_1\vecmu+B_2r_0.
\]
The input is admissible if $B_3-r_1\in\Rq^\times$.  When admissible, let
\[
  Z=(B_3-r_1)^{-1}.
\]
The BB-tran target is
\begin{equation}\label{eq:htran-target}
  T=h_{\mathrm{tran},\vecmu,(r_0,r_1)}(B):=B_0+B_1\vecmu+B_2r_0+Z.
\end{equation}
Equivalently, the witness form is
\begin{equation}\label{eq:htran-witness-form}
  (B_3-r_1)\Had Z=1,
  \qquad
  T=B_0+B_1\vecmu+B_2r_0+Z,
\end{equation}
and the eliminated form is
\begin{equation}\label{eq:htran-eliminated}
  (B_3-r_1)\Had\bigl(T-B_0-B_1\vecmu-B_2r_0\bigr)=1.
\end{equation}
The witness form is used in the issuance and showing relations because it separates the inverse check from
the linear target equation.

\begin{lemma}[Eliminated form under admissibility]\label{lem:hash-cleared-admissible}
Let $\Rq\cong\prod_i F_i$ be the CRT decomposition and write $B_3-r_1=(d_i)_i$,
$T-L_B(\vecmu,r_0)=(z_i)_i$.  If $d_i\ne0$ for all $i$, then
\[
  T=h_{\mathrm{tran},\vecmu,(r_0,r_1)}(B)
  \quad\Longleftrightarrow\quad
  (B_3-r_1)\Had\bigl(T-B_0-B_1\vecmu-B_2r_0\bigr)=1.
\]
\end{lemma}

\begin{proof}
In each CRT slot the defining statement is $z_i=d_i^{-1}$, which is equivalent to $d_i z_i=1$ when
$d_i\ne0$.  Reassembling the CRT slots gives the claim.
\end{proof}

\begin{definition}[Direct target hiding of BB-tran]\label{def:htran-direct-hiding}
We say that the instantiated BB-tran family is directly target-hiding if the following distributions are
computationally indistinguishable for all bounded challenge messages $\vecmu_0,\vecmu_1$.  Sample
$b\getsunif\{0,1\}$, sample $r_0,r_1$ from the protocol randomness distribution conditioned on
$B_3-r_1\in\Rq^\times$, and output
\[
  T=B_0+B_1\vecmu_b+B_2r_0+(B_3-r_1)^{-1}.
\]
No PPT distinguisher, given the public parameters and $T$, can guess $b$ with non-negligible advantage.
\end{definition}

The security proof in Section~\ref{sec:blind-signature} uses this direct target-hiding property as the
privacy input for the rational-hash layer.  It is separate from the non-blind unforgeability import from
DKLW and does not require the target to be computed from the commitment value.

\begin{definition}[vSIS-style hash-and-sign]\label{def:vsis-template}
Let $\TDRH$ be a trapdoor rational hash scheme with parameters $B,\mathcal X,\delta,\BoundSig$.  The
associated signature scheme $\vsisscheme=(\KeyGen,\Sign,\SigVerify)$ is:
\begin{itemize}[leftmargin=1.25em]
  \item $\KeyGen$ runs $(\mA,\trapdoor)\gets\TrapGen$ and outputs $\vk=\mA$, $\sk=\trapdoor$.
  \item $\Sign(\sk,\mu)$ samples $\chi\gets\mathcal X$ until $T=\Eval(B,\mu,\chi)\ne\bot$, samples
  $\vecu\gets\SampPre(\trapdoor,T)$, and outputs $\sigma=(\vecu,\chi)$.
  \item $\SigVerify(\vk,\mu,\sigma)$ for $\sigma=(\vecu,\chi)$ checks $\chi\in\mathcal X$,
  $\Eval(B,\mu,\chi)\ne\bot$, $\|\vecu\|_\infty\le\BoundSig$, and
  $\mA\vecu=\Eval(B,\mu,\chi)$.
\end{itemize}
\end{definition}

\begin{lemma}[Imported non-blind security]\label{lem:vsis-unforgeability}
Assume the trapdoor parameters are admissible, the BB-tran inadmissibility probability is negligible, the
accepted preimages satisfy the norm bound required by~\cite{EPRINT:DKLW25}, and the denominator-bound,
strong-linear-independence, and predicate side conditions required by the BB-tran analysis of~\cite{EPRINT:DKLW25}
hold.  Then the non-blind hash-and-preimage signature layer of Definition~\ref{def:vsis-template} inherits
the corresponding DKLW unforgeability guarantees under the stated strong hinted-vSIS/GenISIS assumptions.
\end{lemma}

\begin{proof}
This is the specialization of~\cite[Theorems~4, 5, 7, 8 and Corollary~1]{EPRINT:DKLW25} to the BB-tran
family, together with the BB-tran-specific GenISIS/IntGenISIS lift of~\cite[Section~5.2]{EPRINT:DKLW25}.
\end{proof}

\subsection{Non-interactive zero-knowledge arguments of knowledge}
\label{subsec:prelim-nizk}
\label{subsec:nizk-ark}
We use the NIZK-ARK syntax of the SmallWood framework instantiated in the random-oracle model
\cite[Definition~3]{EPRINT:FenRiv25b}.

\begin{definition}[NIZKs in the random oracle model]\label{def:nizk-ark}
Let $\mathcal R\subseteq\mathcal X\times\mathcal W$ be a family of relations.  A transparent non-interactive
zero-knowledge argument of knowledge for $\mathcal R$ in the random-oracle model is a pair
$(\ZKProve^{\RO},\ZKVerify^{\RO})$ satisfying correctness, straightline-extractable knowledge soundness,
and zero knowledge.  In particular, zero knowledge means that there exists a simulator
$\Sim^{\RO}$ which, given only the public statement, outputs a proof indistinguishable from a real proof.
\end{definition}

We instantiate these proofs using the SmallWood PACS/PIOP framework compiled with Fiat--Shamir in the
random-oracle model~\cite{EPRINT:FenRiv25b}.

\subsection{Blind signatures}
\label{subsec:blind-sig-defs}
A blind signature scheme allows a holder $\Holder$ to obtain a signature valid under an issuer's verification
key without revealing the signed message to the issuer.  We follow the standard syntax of~\cite{C:Fischlin06}.

\begin{definition}[Blind signature syntax]\label{def:bs-syntax}
A blind signature scheme is a tuple
\[
  \BlindSignature=(\Setup,\IKeyGen,\Issue,\Verify)
\]
where $\Setup$ generates public parameters, $\IKeyGen$ generates issuer keys $(\vk,\sk)$, $\Issue$ is an
interactive protocol between issuer and holder, and $\Verify(\vk,\mu,\sig)$ checks signatures.
\end{definition}

\begin{definition}[Correctness]\label{def:bs-correctness}
A blind signature scheme is correct if an honestly issued signature verifies except with negligible probability.
\end{definition}

\begin{definition}[Blindness]\label{def:bs-blindness}
A blind signature scheme is blind if, for any two messages $\mu_0,\mu_1$, no PPT issuer-side adversary can
distinguish the order of two completed issuance transcripts/signatures except with negligible advantage.
\end{definition}

\begin{definition}[One-more unforgeability]\label{def:bs-one-more}
A blind signature scheme is one-more unforgeable if any PPT adversary that interacts with the honest issuer
in at most $Q$ issuance sessions cannot output $Q+1$ valid signatures on distinct messages except with
negligible probability.
\end{definition}

\paragraph{Rate limiting at the ARC layer.}
We do not introduce a separate stand-alone rate-limiting primitive.  In this paper, rate limiting is part of
the ARC syntax: the holder presents a credential together with a public nonce and a PRF-derived public tag,
and the verifier enforces the policy through state on previously accepted $(\nonce,\prftag)$ pairs.

\subsection{Anonymous Rate-Limiting Credentials}
\label{subsec:arc-security}
Anonymous rate-limited credentials~\cite{ietf-privacypass-arc-protocol-01} extend blind signatures by adding
a showing algorithm that may create multiple presentations of a previously issued credential.

\begin{definition}[ARC scheme syntax]\label{def:arc-syntax}
Let $\nonceSpace$ be a finite nonce set.  An ARC scheme is a tuple
\[
  \ARC=(\Setup,\IKeyGen,\Issue,\Show,\Verify)
\]
where $\Issue$ outputs a credential, $\Show(\credential,\nonce)$ outputs a presentation, and
$\Verify(\vk,\presentation,\algstate)$ checks the presentation and updates verifier state.
\end{definition}

\begin{definition}[ARC soundness]\label{def:arc-soundness}
Let $\ARC$ have nonce space $\nonceSpace$.  For $Q=Q(\lambda)$, set $K=Q|\nonceSpace|+1$.  After at most
$Q$ issuance interactions, no PPT adversary can produce $K$ presentations all accepted by a stateful verifier,
except with negligible probability.
\end{definition}

\begin{definition}[Presentation unlinkability]\label{def:pres-unlink}
An ARC scheme has presentation unlinkability if, for distinct fresh nonces, presentations produced from the
same credential are indistinguishable from presentations produced from independently issued credentials,
except with negligible advantage.
\end{definition}

\subsection{Pseudorandom functions}
\label{subsec:prelim-prf}
\label{subsec:prf}

\begin{definition}[Pseudorandom function]\label{def:prf}
Let $\mathcal K,\mathcal D,\mathcal T$ be sets with $|\mathcal K|$ exponential in $\lambda$.  A function family
$\PRF:\mathcal K\times\mathcal D\to\mathcal T$ is a PRF if oracle access to $\PRF(k,\cdot)$ for uniform
$k\getsunif\mathcal K$ is computationally indistinguishable from oracle access to a uniform random function.
\end{definition}

We instantiate the PRF with a Poseidon2-style permutation~\cite{AFRICACRYPT:GraKhoSch23}.

\begin{definition}[Poseidon2-style permutation]\label{def:poseidon-perm}
Let $q$ be prime, let $d\ge3$ satisfy $\gcd(d,q-1)=1$, and let $t,R_F,R_P\in\mathbb N$ with $R_F$ even.
A Poseidon2-style permutation alternates external rounds, which apply the S-box to all lanes, and internal
rounds, which apply it to one lane, together with public round constants and invertible linear mixing maps.
\end{definition}

\begin{definition}[Poseidon2 PRF]\label{def:poseidon-prf}
For key $\veck\in\Fq^{\ell_{\mathsf{key}}}$ and public input $\vecn\in\Fq^{\ell_{\mathsf{in}}}$, define
\[
  \PRF(\veck,\vecn)=\Tr_{\ell_{\mathsf{tag}}}\bigl(P(\veck\Vert\vecn)+(\veck\Vert\vecn)\bigr),
\]
where $P$ is the Poseidon2-style permutation and $\Tr_{\ell_{\mathsf{tag}}}$ outputs the first
$\ell_{\mathsf{tag}}$ lanes.
\end{definition}
